\lesson{2}{Sep 14 2021 Tue (09:54:31)}{Properties of Rational Exponents}{Unit 1}

\subsubsection{Dividing Rational Exponents}
\label{sub_sub_sec:dividing_rational_exponents}

You already know that when like variables are multiplied, their exponents are
added. But what happens when you divide them? Look at this
property~\ref{prop:quotient_of_a_power_property}

\begin{property}[Quotient of a Power Property]
  \label{prop:quotient_of_a_power_property}

  To divide powers of the same base, subtract the exponents:
\end{property}

\begin{example}
  \label{exm:quotient_of_a_power_property}

  \begin{equation}
    \label{eqt:quotient_of_a_power_property}
    \begin{split}
      \frac{a^{5}}{a^{3}} &= \frac{\cancel{a} \times \cancel{a} \times
                       \cancel{a} \times a \times a}{\cancel{a} \times \cancel{a} \times
                       \cancel{a}}
                    = a^{2} \\
                    &= a^{5 - 3}
    \end{split}
  .\end{equation}
\end{example}

\begin{review}[Finding a Common Denominator]
  \label{rev:finding_a_common_denominator}

  Look at the factors: $\frac{1}{2}$ and $\frac{1}{4}$.

  First, you start by listing the multiples of each denominator:

  \[ 2: 2, \boxed{4}, 6, 8, 10, 12, \ldots \].
  \[ 4: \boxed{4}, 8, 12, 16, 20, 24 \ldots \].

  So, we just found our common denominator: $4$. Since the first fraction:
  $\frac{1}{4}$, already has a denominator of $4$, we can just leave it alone.
  But, the second fraction: $\frac{1}{2}$, needs to have a denominator of $4$.
  So, we just do: $4 \div 2 = 2$. So, we multiply the second fraction by
  $\frac{2}{2}$.

  \[ \frac{1}{2} \times \frac{2}{4} = \frac{2}{4} \].

  Multiplying the numerator and denominator by the same factor is the same as
  multiplying the entire fraction by 1 since $\frac{2}{2} = 1$. By multiplying 
  $\frac{1}{2}$ by $\frac{2}{2}$, the value of the fraction hasn't changed. In
  fact, if you simplify the new fraction we got, which was $\frac{2}{4}$, it
  goes back to the original fraction: $\frac{1}{2}$. But now, it just has a
  common denominator of $4$, which is what we need to add the exponents.
\end{review}

% subsubsection dividing_rational_exponents (end)

\subsubsection{Raising a Power to a Power}
\label{sub_sub_sec:raising_a_power_to_a_power}

Multiplication of powers with the same base involves the addition of exponents.
Division of powers with the same base involves the subtraction of exponents.
What operation is performed on the exponents of $2$ and $3$ to result in the
exponent of $6$? Take a look at this property
property~\ref{prop:raising_a_power_to_a_power}.

\begin{property}[Power of a Power Property]
  \label{prop:power_of_a_power_property}
  To raise a power to a power, multiply the exponents:

  \[ \left(a^{m}\right)^{n} = a^{m \times n} \].
\end{property}

\begin{example}
  \label{exm:power_of_a_power_property}

  \begin{equation}
    \label{eqt:raising_a_power_to_a_power}
    \begin{split}
      \left(c^{\frac{1}{2}}\right)^{\frac{1}{4}} &= c^{\frac{1}{2} \times \frac{1}{4}} \\
                  &= c^{\frac{1}{8}} \\
                  &= \sqrt[8]{c}
    \end{split}
  .\end{equation}
\end{example}

% subsubsection raising_a_power_to_a_power (end)

\subsubsection{Negative Rational Exponents}
\label{sub_sub_sec:negative_rational_exponents}

Let's take a look at the pattern that forms as you look at decreasing powers of
$4$:

\begin{equation}
  \label{eqt:negative_rational_exponents}
  \begin{split}
    4^{4} &= 4 \times 4 \times 4 \times 4 \times 4 = 256 \\
    4^{3} &= 4 \times 4 \times 4 = 64 \\
    4^{2} &= 4 \times 4 = 16 \\
    4^{1} &= 4 = 4 \\
    4^{0} &= 1 \\
    4^{-1} &= \frac{1}{4^{1}} = \frac{1}{4} \\
    4^{-2} &= \frac{1}{4^{2}} = \frac{1}{16} \\
    4^{-3} &= \frac{1}{4^{3}} = \frac{1}{64} \\
    4^{-4} &= \frac{1}{4^{4}} = \frac{1}{256} \\
  \end{split}
.\end{equation}

Recall that the inverse of a positive is a negative, and vise versa. Also, the
inverse of multiplication is division and the inverse of addition is
subtraction. When the exponent is positive, it tells you to multiply the base
the number of times, which is indicated by the exponent.

\[ 4^{3} = 4 \times 4 \times 4 = 64 \].

Inversely, a negative exponent tells you to divide the base by the number of
times indicated by the exponent.

\[ 4^{-3} = 4 \div 4 \div 4 = 0.015625 \].

or

\[ 4^{-3} = \frac{1}{4^{3}} = \frac{1}{64} = 0.015625 \].

\begin{example}
  \label{exm:negative_rational_exponents}

  Simplify the following expression:

  \begin{equation}
    \label{eqt:negative_rational_exponents_simplify_2}
    \begin{split}
      9^{-\frac{1}{2}} &= \frac{1}{9^{\frac{1}{1}}} \\
             &= \frac{1}{\sqrt{9}} \\
             &= \frac{1}{3}
    \end{split}
  .\end{equation}
\end{example}

% subsubsection negative_rational_exponents (end)

\newpage

